
Contrastive self-supervised learning is commonly assumed to amplify spurious correlations: its augmentation-invariance objective makes salient spurious features --- such as background texture --- strongly instance-discriminative. We test this assumption directly with a controlled four-paradigm comparison of frozen ResNet-50 representations on two group-annotated benchmarks. Contrary to the augmentation-invariance prediction, supervised ERM encodes spurious background features more strongly than contrastive MoCo-v3 on Waterbirds (mean ratio difference = 0.0247, Cohen's d = 5.68, ANOVA F = 35.99, p = 2.42 $\times 10^{-7})$, while self-distillation DINO produces ratios statistically indistinguishable from ERM (d = 0.60, p\_bonf = 1.0). On CelebA, the paradigm ranking reverses (MoCo-v3 highest, DINO lowest; ANOVA F = 5.51, p = 0.009), establishing a spurious-attribute-type $\times$ pretraining-objective interaction. When SimCLR is trained from scratch on Waterbirds with background-replacement augmentation, the spurious/task probe accuracy ratio falls by 23.1\% relative to standard augmentation (mean ratio 1.472 vs. 1.132, paired t = 46.66, p = 1.26 $\times 10^{-6}$, d = 32.4), confirming that augmentation design is a functional causal lever for spurious feature encoding. These results indicate that supervised label correlation --- not augmentation-based instance-discriminability --- is the dominant driver of spurious feature encoding in frozen ResNet-50 representations, and that label-agnostic contrastive pretraining partially suppresses spurious features by design.

